\documentclass{article}
\usepackage[T1]{fontenc}
\usepackage{lmodern}
\usepackage{graphicx}
\usepackage{amsmath,amssymb}
\usepackage[a4paper,margin=2cm]{geometry}
\usepackage[absolute,overlay]{textpos}
\setlength{\TPHorizModule}{1mm}
\setlength{\TPVertModule}{1mm}
\usepackage[indonesian]{babel}
\usepackage{hyperref}
\setlength{\emergencystretch}{2em}
% unit: Halaman 33

\begin{document}

% === Halaman 33 (python/native) ===

33

Dekripsi: Bagi panjang cipherteks dengan panjang kunci  (Pada contoh ini, 30 / 6 = 5) .

                Cipherteks: SDNIAIAONSSNLFITTOOIEEGRTMKIMB


Tulis cipherteks secara vertical sepanjang 5 kolom:        




SDNIA



IAONS



SNLFI



TTOOI



EEGRT



MKIMB

Plainteks: (baca secara vertikal)

sistemdanteknologiinformasiitb

sistem dan teknologi informasi itb         (tambahkan spasi)

\end{document}
